\documentclass[10pt,a4paper]{amsart}
\usepackage[margin=19mm]{geometry}
\usepackage{amsmath,amssymb,mathtools}
\usepackage[scr=rsfs,cal=euler]{mathalfa}
\usepackage{xfrac}
\usepackage[unicode,hidelinks,bookmarks=false]{hyperref}
\newcommand{\ie}{i.e.\ }
\newcommand{\cf}{cf.\ }
\newcommand{\resp}{respectively\ }
\newcommand{\Subsection}{\subsection}
\providecommand{\nobreakdash}{\nobreak\hbox{-}}
\setlength{\emergencystretch}{2em}
\multlinegap0pt
\usepackage{bm}
\usepackage{bbm}
\usepackage{dsfont}
\usepackage{xspace}
\usepackage{cancel}
\usepackage{mathtools}
\usepackage{amsthm}
\makeatletter
\@ifundefined{theorem}{\newtheorem{theorem}{Theorem}}{}
\@ifundefined{claim}{\newtheorem{claim}{Claim}}{}
\@ifundefined{claimproof}{\newtheorem{claimproof}{Claimproof}}{}
\@ifundefined{lemma}{\newtheorem{lemma}[theorem]{Lemma}}{}
\makeatother
\makeatletter
\DeclareMathOperator{\dcl}{dcl} 
\newcommand{\forkindep}[1][]{%
  \mathrel{
    \mathop{
      \vcenter{
        \hbox{\oalign{\noalign{\kern-.3ex}\hfil$\vert$\hfil\cr
              \noalign{\kern-.7ex}
              $\smile$\cr\noalign{\kern-.3ex}}}
      }
    }\displaylimits_{#1}
  }
}
\expandafter\ifx\csname step\endcsname\relax
\newenvironment{step}
    {\refstepcounter{step}
     \vspace{7pt}\par\noindent
     \textbf{Step \thestep.}\hspace{0.3mm}}%
    {\hfill $\clubsuit$             %puts a club-symbol at the end of the step
     \vspace{7pt}\par}
\fi
\makeatother
\title[Infinite Cliques in Simple and Stable Graphs]{Infinite Cliques in Simple and Stable Graphs}
\author{Yatir Halevi, Itay Kaplan, Saharon Shelah}
\date{Source: arXiv:2408.05605v1 (CC BY 4.0)}
\begin{document}
\maketitle

\noindent\emph{Source and licence.} Infinite Cliques in Simple and Stable Graphs. Version arXiv:2408.05605v1 (CC BY 4.0), distributed under the Creative Commons Attribution 4.0 International licence, as recorded in the arXiv licence metadata for this exact version and shown on the abstract page https://arxiv.org/abs/2408.05605v1. Full rights notice: licenses/arxiv-2408.05605-CC-BY-4.0.txt.\par\smallskip

\noindent\textit{Editorial scope.} Counted unit: the Lemma whose source label is \texttt{L:technical simple lemma} in the article (source0.tex lines 410--419, source label \texttt{L:technical simple lemma}), delivered together with the article's own complete original proof of it (source0.tex lines 420--448, one \texttt{proof} environment of 29 lines). No mathematics is written, repaired or re-derived: statement and proof are the article's own text line for line. Required context retained uncounted: none. Every definition, display and label of those lines is retained unchanged. Omitted from this package and not redistributed: 1--409, 449--1330. Nothing inside a retained excerpt is omitted or altered: every retained body is delivered exactly as the article's lines are written, so no omission receipt of any kind accompanies this package. Citations: the retained excerpts carry no citation command at all, so no bibliography entry of the article is transcribed and no reference is redistributed. Reference adaptation: none was needed and none was made; every reference inside the delivered unit resolves inside it, so no unresolved reference or citation is produced. Printed slips kept literal: no printed word, symbol, hypothesis or number of the counted statement or of its proof was altered. Layout and class: the article's own class is \texttt{amsart} with packages including amssymb, latexsym, inputenc, amsmath, amsthm, amsfonts, mathrsfs, hyperref, cleveref, enumerate, units, xy, graphicx, url; the wrapper is the lane's pinned \texttt{amsart} editorial wrapper at 10pt on A4, which restates the theorem-like environments the retained text needs and adds the article's own macro definitions that the retained text needs (\texttt{\textbackslash dcl}, \texttt{\textbackslash forkindep}), with extra wrapper packages (bm, bbm, dsfont, xspace, cancel, mathtools). The delivered numbering is the unit's own. Assets: no retained span contains a figure environment, a graphics-inclusion command, a PDF-page inclusion, a table or a TikZ picture; no image file is copied into this package, the package declares no asset dependency and no image file is redistributed with it. Licence: version arXiv:2408.05605v1 of this article is distributed under the Creative Commons Attribution 4.0 International licence, as recorded in the arXiv licence metadata for this exact version and shown on the abstract page https://arxiv.org/abs/2408.05605v1. A full rights notice is delivered at licenses/arxiv-2408.05605-CC-BY-4.0.txt. Characters: every retained excerpt is pure ASCII, so the delivered engine, XeLaTeX with its default Latin Modern text font, reports no missing character.
\begin{lemma}\label{L:technical simple lemma}
Let $M$ be some structure, in a language $L$, and assume that $T=\mathrm{Th}(M)$ is simple. Let $M'$ be some expansion of $M$ to a language $L'\supseteq L$. Let $<$ be a linear order on the universe of $M$ and $\alpha,\beta\in M$ elements satisfying that
\begin{enumerate}
\item $(\dcl_{L'}(\alpha),<)$ and $(\dcl_{L'}(\beta),<)$ are well-orders and
\item  for any $\gamma\in \dcl_{L'}(\alpha)\cup \dcl_{L'}(\beta)$, $\gamma \forkindep[\dcl_{L'}(\gamma)\cap \{\varepsilon\in M:\varepsilon<\gamma\}]\{\varepsilon\in M:\varepsilon<\gamma\}$.
\end{enumerate}


Then  \[\dcl_{L'}(\alpha)\forkindep[\dcl_{L'}(\alpha)\cap \dcl_{L'}(\beta)]\dcl_{L'}(\beta).\]
\end{lemma}

\begin{proof}
Let $\alpha,\beta\in M$ be elements as in the statement, and let $\mathbf{a}^\alpha=\dcl_{L'}(\alpha)$ and $\mathbf{a}^\beta=\dcl_{L'}(\beta)$.
Set $\Omega=\mathbf{a}^\alpha\cup\mathbf{a}^\beta$ and for any $\gamma\in \Omega$ let $\Omega_{<\gamma}=\{\varepsilon \in \Omega: \varepsilon<\gamma\}$. Since $(\Omega,<)$ is a finite union of well ordered sets, it is also well ordered. 

\begin{claim}\label{C:induction}
For $\gamma\in \Omega$, if \[\mathbf{a}^\alpha\cap \Omega_{<\gamma} \forkindep[\mathbf{a}^\alpha\cap \mathbf{a}^\beta\cap \Omega_{<\gamma}]\mathbf{a}^\beta\cap\ \Omega_{<\gamma}\]
then
\[\mathbf{a}^\alpha\cap \Omega_{\leq \gamma} \forkindep[\mathbf{a}^\alpha\cap \mathbf{a}^\beta\cap \Omega_{\leq \gamma}]\mathbf{a}^\beta\cap\ \Omega_{\leq \gamma}.\]
\end{claim}
\begin{claimproof}
By symmetry, we deal with the case $\gamma\in \mathbf{a}^\alpha$. We first prove that
\begin{equation}\label{eq:indep}
\mathbf{a}^\alpha\cap \Omega_{\leq \gamma}\forkindep[\mathbf{a}^\alpha\cap \mathbf{a}^\beta\cap \Omega_{<\gamma}] \mathbf{a}^\beta\cap \Omega_{<\gamma}.
\end{equation}

By hypothesis (2), $\gamma\forkindep[\dcl_{L'}(\gamma)\cap \{\varepsilon:\varepsilon<\gamma\}]\{\varepsilon:\varepsilon<\gamma\}$.
Since \[\dcl_{L'}(\gamma)\cap\{\varepsilon:\varepsilon<\gamma\}\subseteq \mathbf{a}^\alpha\cap \Omega_{<\gamma}\subseteq \{\varepsilon:\varepsilon<\gamma\},\] we get that $\gamma\forkindep[\mathbf{a}^\alpha\cap \Omega_{<\gamma}]\{\varepsilon:\varepsilon<\gamma\}$, so $\gamma\forkindep[\mathbf{a}^\alpha\cap \Omega_{<\gamma}] \mathbf{a}^\beta\cap \Omega_{<\gamma}$ and \[\mathbf{a}^\alpha\cap\Omega_{\leq \gamma}\forkindep[\mathbf{a}^\alpha\cap \Omega_{<\gamma}] \mathbf{a}^\beta\cap \Omega_{<\gamma}.\] By assumption and transitivity, we conclude (\ref{eq:indep}).

If $\gamma\notin \mathbf{a}^\beta$ then we are done; so assume that $\gamma\in \mathbf{a}^\beta$ as well. In this case, by properties of forking we get that $\mathbf{a}^\alpha\cap  \Omega_{\leq \gamma}\forkindep[\mathbf{a}^\alpha\cap \mathbf{a}^\beta\cap \Omega_{\leq \gamma}]\mathbf{a}^\beta\cap\Omega_{\leq \gamma}$, which is what we wanted to prove.
\end{claimproof}
The proof now follows by induction: the successor step is Claim \ref{C:induction} and limit case follows since forking is witnessed by a formula.
%
%If $\gamma\notin \mathbf{a}^\alpha\cup \mathbf{a}^\beta$ then $\mathbf{a}^\alpha\cap (\gamma+1)=\mathbf{a}^\alpha\cap (\gamma\cup \{\gamma\})=\mathbf{a}^\alpha\cap \gamma$, and similarly for $\mathbf{a}^\beta$, so there is nothing to show. 
%
%Otherwise there are two options, either $\gamma\in \mathbf{a}^\alpha$ or $\gamma\in \mathbf{a}^\beta\setminus \mathbf{a}^\alpha$ or $\gamma\in \mathbf{a}^\alpha\cap \mathbf{a}^\beta$.
%
%Assume that $\gamma\in \mathbf{a}^\alpha\setminus \mathbf{a}^\beta$. By assumption we need to show that 
%Assume that $\gamma\in \mathbf{a}^\alpha\cap \mathbf{a}^\beta$. We want to prove that $(\mathbf{a}^\alpha\cap \gamma)\cup\{\gamma\}\forkindep[(\mathbf{a}^\alpha\cap \mathbf{a}^\beta\cap \gamma)] (\mathbf{a}^\beta\cap \gamma)$. As $\dcl_{L'}(\{\gamma\})\cap \gamma\subseteq \mathbf{a}^\alpha\cap \mathbf{a}^\beta\cap \gamma$ and $\{\gamma\}\forkindep[\dcl_{L'}(\{\gamma\})\cap\gamma]\gamma$ we get that $\{\gamma\}\forkindep[\mathbf{a}^\alpha\cap \mathbf{a}^\beta\cap \gamma] \gamma$ and thus $\{\gamma\}\forkindep[(\mathbf{a}^\alpha\cap \mathbf{a}^\beta\cap\gamma)\cup (\mathbf{a}^\alpha\cap \gamma)] (\mathbf{a}^\beta\cap \gamma)$, so  $\{\gamma\}\forkindep[\mathbf{a}^\alpha\cap \gamma] (\mathbf{a}^\beta\cap \gamma)$. Consequently, $(\mathbf{a}^\alpha\cap \gamma)\cup\{\gamma\}\forkindep[\mathbf{a}^\alpha\cap \gamma] (\mathbf{a}^\beta\cap \gamma)$. As above, it follows by the induction hypothesis and transitivity that $(\mathbf{a}^\alpha\cap \gamma)\cup\{\gamma\}\forkindep[(\mathbf{a}^\alpha\cap \mathbf{a}^\beta\cap\gamma)] (\mathbf{a}^\beta\cap \gamma)$.
\end{proof}

\end{document}
